\section{Complexity: the cross-over is where it was predicted, the advantage is not the size that was claimed}
\label{sec:complexity}

Three claims are usually made together about the efficiency of state space
vision backbones: that a cross-over exists beyond which they are cheaper than
self-attention, that the advantage past that point is large, and that it
translates into throughput and memory at high resolution. They are separable,
and measurement separates them.

\begin{figure}[tbp]
  \centering
  \includegraphics[width=\linewidth]{fig1_complexity.pdf}
  \caption{Measured FLOPs and throughput against input resolution for the
  three whole models. In the FLOPs panel Vim-S sits above DeiT-S at
  $224^2$ and $384^2$ and below it from $512^2$ on: the cross-over falls in
  the interval the analysis predicts. Cells labelled \emph{n/a} in the
  throughput panel were not measured---the maximum-batch search exhausted
  memory---and the connecting line spans them.}
  \label{fig:complexity}
\end{figure}

\begin{table}[t]
\centering
\caption{Measured resolution sweep in the pinned environment. \textemdash{} marks a cell where the throughput search hit CUDA OOM; those runs still completed every other measurement, and the reason is recorded in the \texttt{error} column of \texttt{results/e1/sweep.csv}.}
\label{tab:e1}
\begin{tabular}{llrrrrr}
\toprule
Model & Res. & Params (M) & FLOPs (G) & Latency (ms) & Peak alloc (MiB) & img/s \\
\midrule
DeiT-S & 224 & 22.05 & 4.61 & 3.16 & 97.2 & 1038.0 \\
DeiT-S & 384 & 22.20 & 15.52 & 5.46 & 105.3 & 282.0 \\
DeiT-S & 512 & 22.37 & 31.80 & 10.95 & 115.6 & 130.0 \\
DeiT-S & 768 & 22.86 & 98.70 & 31.32 & 144.2 & 36.9 \\
DeiT-S & 1024 & 23.55 & 243.09 & 78.89 & 182.6 & 13.9 \\
CMT-S & 224 & 26.26 & 4.09 & 18.74 & 243.9 & 456.6 \\
CMT-S & 384 & 28.46 & 13.00 & 20.03 & 265.5 & 133.8 \\
CMT-S & 512 & 33.83 & 25.20 & 21.46 & 301.1 & \textemdash \\
CMT-S & 768 & 65.78 & 70.09 & 45.67 & 466.8 & 22.8 \\
CMT-S & 1024 & 151.80 & 157.96 & 101.73 & 1176.9 & \textemdash \\
Vim-S & 224 & 25.80 & 5.78 & 18.13 & 116.1 & 373.6 \\
Vim-S & 384 & 25.94 & 16.93 & 18.41 & 132.4 & 100.3 \\
Vim-S & 512 & 26.11 & 30.08 & 19.47 & 149.7 & \textemdash \\
Vim-S & 768 & 26.61 & 67.64 & 37.69 & 199.4 & 25.7 \\
Vim-S & 1024 & 27.29 & 120.22 & 59.49 & 267.0 & 17.7 \\
\bottomrule
\end{tabular}
\end{table}


\subsection{The cross-over point is confirmed}

The analytical argument places the cross-over where the quadratic term of
self-attention overtakes its linear term, at a token count of roughly twice
the channel width. For the widths used here that is on the order of several
hundred tokens, which falls between the $384^2$ and $512^2$ inputs of our
sweep.

Table~\ref{tab:e1} puts it there. At $384^2$ Vim-S costs more total FLOPs
than DeiT-S; at $512^2$ it costs less, and the gap widens with resolution.
The prediction was made about a term in an expression and it survives being
checked against whole measured models, which is not guaranteed---the
expression describes one block, and a model is not one block.

\subsection{The size of the advantage does not survive}

The comparison table that accompanies this claim in the literature we are
auditing reports a ratio of more than fifty to one at a large token count.
Measured end to end, at $1024^2$, Vim-S costs \VimFlopsHigh~GFLOPs against
DeiT-S at \DeitFlopsHigh~GFLOPs. That is a factor of two.

The discrepancy is not noise, and it is not a different model. It comes from
an asymmetry in what was counted, which we take up in
Section~\ref{sec:discussion}: the self-attention figure includes the
projection terms that scale with the square of the channel width, and the
state space figure counts only the scan. A Mamba block also has input,
output, and parameter projections of that order, and once they are included
the fifty-fold gap becomes a two-fold one.

There is a second reason to distrust the common way of producing these
numbers, and it points the same direction. Automatic FLOP counters~\citep{fvcore} attribute
zero cost to fused kernels they cannot trace into. In this repository that
happened twice: a fused attention kernel swallowed the attention matrix
multiplications of DeiT, and a fused selective-scan kernel swallowed the
convolution, the projections, and the scan of Vim. Both models are
undercounted by the naive procedure, but the state space model is undercounted
far more, because more of its arithmetic lives inside one fused call. A
measurement taken without that care makes SSMs look cheaper than they are.
We therefore un-fuse the kernels while counting, which is legitimate because
fusion changes the kernel boundary and not the arithmetic, and we never
measure latency, memory, or throughput in that mode, where fusion is the
performance.

\subsection{The throughput and memory claims do not reproduce here}

The efficiency figures most often quoted for Vim---roughly a threefold
throughput advantage and a large reduction in memory at high
resolution---were reported at a larger input than our sweep reaches and on
different hardware. We do not reproduce them at $1024^2$ on an
\GpuName{}. Vim-S runs at \VimImgsHigh~images per second against DeiT-S at
\DeitImgsHigh{}, a ratio well below the reported one, and its peak allocated
memory is \VimPeakHigh~MiB against \DeitPeakHigh~MiB---higher, not lower.
The same ordering appears in how many images fit at once:
\VimMaxBatchHigh{} for Vim-S against \DeitMaxBatchHigh{} for DeiT-S.

This is a failure to reproduce under stated conditions, not a refutation.
The conditions differ in resolution, in device, and in memory budget, and any
of those could account for the gap. What does not depend on the comparison
with the original report is the direction we measure: at this resolution, on
this device, at batch size one, the state space model is the one that runs
out of memory first. The intuition that a linear-time sequence operator is
also the memory-frugal choice does not hold here, because the scan's
intermediate activations have to be kept.

\subsection{A cost that the claims do not mention}

CMT-S is not resolution-scalable in parameter count. It grows from
\CmtParamsLow~M at $224^2$ to \CmtParamsHigh~M at $1024^2$, because its
relative position bias tables are sized by the token grid. Over the same
range DeiT-S and Vim-S grow by about one and one and a half million
parameters respectively. Any high-resolution comparison that treats CMT as a
fixed-size model is comparing something else.

\subsection{What not to quote from this table}

Latency at $224^2$ and $384^2$ should not be cited as a single number. Within
one process the three repeated measurement blocks agree to within a factor of
\LatencySpreadHigh{}, but the same cell measured in two independent runs of
the whole sweep differed by a factor of \LatencyBetweenRunSpread{}. The
uncertainty lives at the process boundary, which repetition inside a process
cannot see. At $768^2$ and above both spreads fall within a few percent and
single values are safe to quote. The two independent sweeps are both in the
repository's history, so this spread can be recomputed rather than taken on
trust.
